\documentclass[a4paper,12pt]{book}
\usepackage[vietnamese,english]{babel}
\usepackage{lmodern}
\usepackage{booktabs}
\usepackage{graphicx}
\usepackage{titling}
\usepackage[Bjornstrup]{fncychap}
\usepackage{hyperref}
\usepackage[titletoc]{appendix}
\usepackage[numbib]{tocbibind}
\usepackage[acronym,numberedsection=autolabel,section,toc]{glossaries}

\author{Nguyễn Gia Phong}
\newcommand{\doi}[1]{doi:\href{https://doi.org/#1}{#1}}

\makeglossaries
\newglossaryentry{content}{name=content,description={a file or a directory}}

\newacronym{usth}{USTH}{University of Science and Technology of Hanoi}
\newacronym{api}{API}{application programming interface}
\newacronym{id}{ID}{identifier}
\newacronym{cid}{CID}{\gls{content} \gls{id}}

\begin{document}
\frontmatter

% Template: https://www.latextemplates.com/template/academic-title-page
\pagestyle{empty}
\begin{center}\large
  \textsc{\Large\selectlanguage{vietnamese}\begin{tabular}{c}
    University of Science\\
    and Technology of Hanoi
  \end{tabular}}
  \vfill\vfill
  \textsc{Bachelor Thesis}
  \vfill
  \textbf{\LARGE\begin{tabular}{c}
    \toprule
    % TODO: Convolute the title
    Development of Data Lake Core\\
    with Append-Only Storages\\
    and Query Polymorphism\\
    \bottomrule
  \end{tabular}}
  \vfill
  \begin{minipage}[t]{0.4\textwidth}
    \begin{flushleft}
      \textit{Author}\\
      \selectlanguage{vietnamese}
      \theauthor
    \end{flushleft}
  \end{minipage}
  \hspace{0.12\textwidth}
  \begin{minipage}[t]{0.4\textwidth}
    \begin{flushright}
      \textit{Supervisor}\\
      \selectlanguage{vietnamese}
      Trần Giang Sơn, PhD
    \end{flushright}
  \end{minipage}
  \vfill\vfill\vfill
  \today
  \vfill
  \includegraphics[width=0.3\textwidth]{usth-logo.pdf}
  \vfill
\end{center}
\cleardoublepage

\setcounter{page}{1}
\pagestyle{plain}
\tableofcontents

\chapter{Declaration}
I declare that I have composed this thesis in its entirety, as the result
of my own work, unless explicitly indicated otherwise via referencing.
The presented work has not been submitted for any previous application
for a degree or professional qualification.

{\selectlanguage{vietnamese}\let\clearpage\relax\chapter*{Lời Cam Đoan}
\vspace{-4em}Tôi xin cam đoan rằng tôi đã soạn toàn bộ luận án này, hoàn toàn
từ kết quả công việc của chính bản thân tôi, trừ khi được chỉ rõ qua trích dẫn.
Công việc được trình bày chưa bao giờ được nộp cho việc cấp học vị
hay chứng chỉ trình độ chuyên môn nào trước đây.}

%\chapter{Preface}
%\chapter{Abstract}

\mainmatter
\pagestyle{headings}
\chapter{Introduction}
\section{Motivation}
Many researchers at \gls{usth} operate with data on a regular basis
and often a dataset is studied by multiple researchers from different
departments and points in time.  Currently the data are organized manually,
even on the laboratories' storages, which is prone to duplication
and makes data discovery difficult.

A data lake shared among the university's researchers, professors and students
will not only save resources but also improve productivity and promote
interdisciplinary collaborations.  With \gls{usth}'s goal of growing to be
an excellent research university in {\selectlanguage{vietnamese}Việt Nam}
and in the region~\cite{usth}, building such data lake can be an essential task.

\section{Background}
A \emph{data lake} is a massive repository of multiple types of data
in their raw format at scale for a low cost~\cite{lake}.
The data's schema (structure) is defined on read to minimize data modeling
and integration costs~\cite{lake}.

For the ease and efficiency of scaling, a microservice architecture
could be a good choice.  By arranging the data lake as a collection
of loosely-coupled services, it becomes possible to scale individual services
individually~\cite{micro}.  In this architecture, the \emph{core} microservice
is defined as the innermost component, which communicates directly
with the storages.  The core shall provide an \gls{api} for other components
to upload, query and extract data.

\emph{Append-only storages} only allow new data to be appended, whilst ensure
the immutability of existing data.  As immutable data are thread-safe,
they reduce the complexity of the concurrency model, making it easier
to comprehend and reason about~\cite{pure}.  This is particularly useful
in large distributed systems with multiple moving parts.

Since the data are immutable, each \emph{\gls{content}} can be given
an \gls{id}, i.e.~a \gls{cid}.  For end-users, we also introduce
a higher level concept: \emph{dataset}, composing of not only the content
but also relevant metadata for indexing.  Like contents, datasets can also
immutable, with changes written as a new revision linking to the previous one.

Append-only storages' operations boil down to two kinds: appending and reading.
For the latter, sometimes the data are not wanted in their entirety,
but filtered and accumulated.  While data of different types usually requires
different tools and libraries to query upon, the core \gls{api} should be
providing one single query language for all data types, plus their metadata.
In this report, such usage is referred to as \emph{query polymorphism}.

\section{Objectives}
The work presented here was done as part of a three-month internship
in collaboration with several other students
at \gls{usth} ICTLab\footnote{\url{https://ictlab.usth.edu.vn}} to build
a data lake for a better management of the university's data.  The internship
focused on the lake's core microservice, which abstracts underlying
persistent layers and perform relevant metadata transformation and discovery.
It should provide an internal interface to other components for data ingestion,
(primitive) query and extraction, as well as carrying out tasks for enhancing
the discoverability and usability of the aforementioned datasets.

\section{Expected Outcomes}
The intended deliverables of the three-month internship are listed as follows:
\begin{itemize}
  \item Requirement analysis of the data lake core
  \item Data lake core's architecture and design
  \item Core \gls{api} design and specification
  \item Implementation and integration with other components
\end{itemize}

\chapter{Methodology}
\section{Requirement Analysis}
As previously introduced, the most basic functions of the data lake core
are content uploading and downloading, along with datasets addition
and querying.  A more advanced (and rather powerful) use case is
content extraction, which allows one to fetch only the interested part
of the content, e.g.~extracting rows matching a certain predicate
from (semi-)structured data.  Together with logging, the core's use cases
are summarized in figure \ref{uc}.

\begin{figure}
  \includegraphics[width=\textwidth]{diagrams/uc.eps}
  \caption{Use-case diagram}
  \label{uc}
\end{figure}

\subsection{Upload content}
\subsection{Add dataset}
\subsection{Find datasets}
\subsection{Download content}
\subsection{Extract content}
\subsection{Gather logs}

\section{Architecture}
\begin{figure}
  \includegraphics[width=\textwidth]{diagrams/arch.eps}
\end{figure}

\section{Design}
\subsection{Technology Choices}

\subsection{Interface}
\begin{figure}
  \includegraphics[width=\textwidth]{diagrams/api.eps}
  \caption{Core HTTP API endpoints in a common order of access}
  \label{db}
\end{figure}

\subsection{Database Schema}
\begin{figure}
  \includegraphics[width=\textwidth]{diagrams/db.eps}
  \caption{Database schema}
  \label{db}
\end{figure}

\subsection{Query Abstract Syntax Tree}

\section{Implementation}

\section{Quality Assurance}

\chapter{Results and Discussion}
\section{Results}
\section{Discussion}

\chapter{Conclusion}

\begin{appendices}
\chapter{Acknowledgment}

\begin{thebibliography}{69}
  \bibitem{usth} \emph{Mission and Vision}.  \acrlong{usth}.
    \url{https://usth.edu.vn/en/abouts/Mission-et-Vision.html}.
  \bibitem{lake} Huang Fang.
    ``Managing Data Lakes in Big Data Era: What's a data lake
    and why has it became popular in data management ecosystem''.
    \emph{2015 IEEE International Conference on Cyber Technology
    in Automation, Control, and Intelligent Systems}, pp.~820--824.
    IEEE, 2015.  \doi{10.1109/CYBER.2015.7288049}
  \bibitem{micro} Chris Richardson.
    ``1.4.1 Scale cube and microservices''.  \emph{Microservice Patterns}.
    Manning Publications, 2018.  ISBN 9781617294549.
  \bibitem{pure} Brian Göetz, Tim Peierls, Joshua Bloch,
    Joseph Bowbeer and David Holmes.
    ``3.4. Immutability''.  \emph{Java Concurrency in Practice}.
    Addison Wesley Professional, 2006.  ISBN 9780321349606.
\end{thebibliography}

\chapter{Terms and Acronyms}
\printglossaries
\end{appendices}
\end{document}
